\title{AlphaFold Protein Structure Database in 2024: providing structure coverage for over 214 million protein sequences}
\begin{document}
\maketitle

\begin{abstract}
The AlphaFold Database Protein Str uct ure Database (AlphaFold DB, https://alphafold.ebi.ac.uk ) has significantly impacted str uct ural biology by amassing o v er 214 million predicted protein str uct ures, expanding from the initial 300k str uct ures released in 2021. Enabled by the groundbreaking AlphaFold2 artificial intelligence (AI) system, the predictions archived in AlphaFold DB have been integrated into primary data resources such as PDB, UniProt, Ensembl, InterPro and MobiDB. Our manuscript details subsequent enhancements in dat a archiving , co v ering successiv e releases encompassing model organisms, global health proteomes, Swiss-Prot integration, and a host of curated protein datasets. We detail the data access mechanisms of AlphaFold DB, from direct file access via FTP to advanced queries using Google Cloud Public Datasets and the programmatic access endpoints of the database. We also discuss the impro v ements and services added since its initial release, including enhancements to the Predicted Aligned Error vie w er, customisation options f or the 3D vie w er, and impro v ements in the search engine of AlphaFold DB. Gr aphical abstr act
\end{abstract}

\section{Recovered Text}
Nucleic Acids Research , 2024, 52 , D368–D375
https://doi.org/10.1093/nar/gkad1011
Advance access publication date: 2 November 2023
Database issue
AlphaF old Prot ein Structure D atabase in 2024: pr oviding
structure co v er ag e f or o v er 214 million prot ein sequences
Mihaly Varadi
1 , Damian Bertoni 1 , Paulyna Magana 1 , Urmila Paramval 1 , Ivanna Pidruchna 1 ,
Malarvizhi Radhakrishnan 1 , Maxim Tsenkov 1 , Sreenath Nair 1 , Milot Mirdita 2 , Jingi Yeo 2 ,
Oleg Kovalevskiy 3 , Kathryn Tuny asuvunakool 3 , A g ata La y don 3 , Augustin Žídek 3 ,
Hamish Tomlinson 3 , Dhavanthi Hariharan 3 , Josh Abrahamson 3 , Tim Green 3 , John Jumper 3 ,
Ew an Birne y
1 , Mar tin St einegg er 2 , Demis Hassabis 3 , * and Sameer Velankar
1 , *
1 European Molecular Biology Laboratory, European Bioinformatics Institute, Hinxton, UK
2 School of Biological Sciences, Seoul National University, Seoul, South Korea
3 Google DeepMind, London, UK
* To whom correspondence should be addressed. Email: sameer@ebi.ac.uk
Correspondence may also be addressed to Demis Hassabis. Email: dhcontact@deepmind.com
Abstract
The AlphaFold Database Protein Str uct ure Database (AlphaFold DB, https://alphafold.ebi.ac.uk ) has significantly impacted str uct ural biology by
amassing o v er 214 million predicted protein str uct ures, expanding from the initial 300k str uct ures released in 2021. Enabled by the groundbreaking
AlphaFold2 artificial intelligence (AI) system, the predictions archived in AlphaFold DB have been integrated into primary data resources such
as PDB, UniProt, Ensembl, InterPro and MobiDB. Our manuscript details subsequent enhancements in dat a archiving , co v ering successiv e
releases encompassing model organisms, global health proteomes, Swiss-Prot integration, and a host of curated protein datasets. We detail
the data access mechanisms of AlphaFold DB, from direct file access via FTP to advanced queries using Google Cloud Public Datasets and
the programmatic access endpoints of the database. We also discuss the impro v ements and services added since its initial release, including
enhancements to the Predicted Aligned Error vie w er, customisation options f or the 3D vie w er, and impro v ements in the search engine of
AlphaFold DB.
Gr aphical abstr act
Introduction
In the past few years, the landscape of protein structure pre-
diction has evolved significantly due to the advent of next-
generation tools such as AlphaFold ( 1 ), RoseTTAFold ( 2 ), and
OpenFold ( 3 ), among others ( 4 ). The development of these
tools was enabled by decades of research in protein sequences
and structures and underscored the importance of open data
and fundamental data resources like the Protein Data Bank
(PDB) ( 5 ) and the Universal Protein Resource (UniProt) ( 6 ).
The new generation of predicted protein structure mod-
els has demonstrated remarkable accuracy, which could mit-
igate the continually widening gap between known protein
sequences and experimentally determined protein structures
( 7 ). Accessing accurate protein structures paves the way for
Received: September 29, 2023. Revised: October 13, 2023. Editorial Decision: October 16, 2023. Accepted: October 18, 2023
© The Author(s) 2023. Published by Oxford University Press on behalf of Nucleic Acids Research.
This is an Open Access article distributed under the terms of the Creative Commons Attribution License (http: // creativecommons.org / licenses / by / 4.0 / ),
which permits unrestricted reuse, distribution, and reproduction in any medium, provided the original work is properly cited.
Downloaded from https://academic.oup.com/nar/article/52/D1/D368/7337620 by guest on 29 April 2026
Nucleic Acids Research , 2024, Vol. 52, Database issue
D 369
an enhanced understanding of protein function, providing re-
searchers with the tools necessary for modulating these pro-
teins or engineering new ones ( 8–10 ). As a result, numerous
areas within the life sciences have experienced considerable
impact due to the availability of accurately predicted pro-
tein structures. The primary domains notably affected include
structure determination, structure-based drug discovery, and
structural bioinformatics ( 11–13 ).
In the wake of the latest breakthroughs in prediction soft-
ware, new databases such as the AlphaFold Protein Struc-
ture Database and the ESM Atlas ( 14 ) have emerged. With
the availability of performant prediction software, a pertinent
question arises: Why is there a need for databases for pre-
dicted structures when researchers can run the software on
their proteins of interest?
While many new-generation protein structure predictors
are readily accessible for the scientific community to run on
arbitrary input protein sequences, a significant barrier exists
for researchers less acquainted with using scientific software.
Moreover, in large-scale bioinformatics analysis like compar-
ative or functional analysis, the necessity to predict a mas-
sive number of protein structures becomes computationally
expensive and redundant, contributing to a significant carbon
footprint. Even in the case of single protein predictions, look-
ing up pre-generated structures takes seconds compared to
potentially hours of computation time. The absence of pre-
generated structures also obstructs the integration of these
valuable predictions into core data resources like UniProt ( 6 ),
InterPro ( 15 ), Ensembl ( 16 ) or the PDBe—Knowledge Base
( 17 ).
This paper presents the data updates and functionality im-
provements in the AlphaFold Protein Structure Database, a
collaborative project between EMBL-EBI and Google Deep-
Mind, since its initial launch in July 2021. We outline how the
coverage of sequence space has improved from the initial re-
lease, with the structure count increasing from approximately
300k to over 214 million (Figure 1 ). Additionally, we pro-
vide insights into the changes and additions to the metadata
and the format of confidence metrics such as the Predicted
Aligned Error (PAE). Updates on accessing the data and the
data types available through File Transfer Protocol (FTP), our
Application Programming Interface (API), and bulk download
options will be discussed. Lastly, we give an overview of all the
improvements and new functionalities introduced on the Al-
phaFold DB website and give an overview of possible future
directions.
Implementation
Updating the AlphaFold DB is a multifaceted process, encom-
passing multiple stages of data management. This process in-
cludes generating a vast array of protein structure predictions,
organising these predictions in a structured and searchable
format, and ensuring straightforward and efficient data ac-
cess for users. The ultimate result of this process is a compre-
hensive and user-oriented experience, facilitating cutting-edge
research across diverse fields within the life sciences.
Data generation
The data generation process for the AlphaFold DB is carried
out by Google DeepMind, with all predictions stored in PDB,
mmCIF and binaryCIF formats, along with their correspond-
ing metadata in JSON format. The generated mmCIF files ad-
here to the modelCIF format ( 18 ).
AlphaFold yields a per-residue estimate of its confidence,
known as pLDDT, ranging from 0 to 100, indicating the tool’s
predicted score on the lDDT-C α metric ( 19 ). The residue-wise
pLDDT scores are stored in the B-factor fields of PDB files
and under the ‘\_ma\_qa\_metric\_local’ category in mmCIF files.
Regions with pLDDT greater than 90 are generally modelled
with high accuracy, making them suitable for high accuracy-
dependent applications such as characterising binding sites.
Those with pLDDT between 70 and 90 are generally well-
modelled, representing a reliable backbone prediction. On the
other hand, regions with pLDDT between 50 and 70 have
lower confidence and should be used cautiously . Finally , re-
gions with pLDDT scores below 50 often exhibit a spaghetti-
like appearance in 3D view, indicating probable disordered
regions. Structured domains with many inter-residue contacts
are generally more reliable than extended linkers or isolated
long helices. Unphysical bond lengths and clashes usually do
not appear in the confident structured regions, and any re-
gion with several of these should be disregarded. Regardless
of the absolute pLDDT score, the PDB and mmCIF files pro-
vide coordinates for all regions, and it is incumbent upon the
user to interpret the model prudently, in line with the provided
guidance.
In addition to the predicted atomic coordinates and the
pLDDT scores, AlphaFold generates a ‘Predicted Aligned Er-
ror’ (PAE) output, representing the predicted error between
relative positions of residue pairs. PAE is a measure indicating
confidence in the relative positions of larger structural units of
proteins, like domains. The PAE can be used to assess the spa-
tial location of protein domains by looking at the values for
residue pairs between different domains. The raw data with
PAE for all residue pairs can be downloaded as a JSON file.
However, parsing the JSON file requires Python or another
programming language for analysis or visualisation.
In 2022, the JSON file format was updated to
a more compact representation. It now consists of
a ‘predicted\_aligned\_error’ field instead of the 1D
‘distances’ field in the earlier representation and a
‘ max\_predicted\_aligned\_error ’ field indicating the maxi-
mum possible value of PAE.
Data archiving
Data archiving for AlphaFold DB began with an initial release
in July 2021, housing over 360 000 structures for 20 model
organism proteomes with sequences derived from the ‘one se-
quence per gene’ reference proteomes provided in UniProt re-
lease 2021\_02. In December 2021, most of the reviewed se-
quences in UniProt, i.e. the Swiss-Prot dataset, were incorpo-
rated from the UniProt release 2021\_04. In January 2022,
proteomes relevant to global health, derived from priority
lists by the World Health Organization, were added, utilising
sequences from UniProt release 2021\_04 ‘one sequence per
gene’ reference proteomes. By July 2022, most of the remain-
ing sequences from UniProt release 2021\_04 were included,
featuring an additional TAR file on the AFDB download page,
EMBL-EBI’s FTP and Google Cloud Datasets, containing pre-
dictions in MANE select ( 20 ).
A November 2022 update rectified structures affected by
a temporary numerical bug presented in the July release. This
bug led to low accuracy predictions with correspondingly low
Downloaded from https://academic.oup.com/nar/article/52/D1/D368/7337620 by guest on 29 April 2026
D 370
Nucleic Acids Research , 2024, Vol. 52, Database issue
Figure 1. The expansion of AlphaFold DB. The AlphaFold Protein Str uct ure Database increased in size through consecutive releases. As of September
2023, it archives over 214 million predicted protein str uct ures.
pLDDT for ∼4\% of the total structure predictions in the
database. As part of this update, the coordinates for affected
structures were updated (old coordinate files remain accessi-
ble as v3 files), and minor metadata adjustments were made
in the mmCIF files for the remaining structures. We docu-
ment every data version update in our changelog at https:
// ftp.ebi.ac.uk/ pub/ databases/ alphafold/ CHANGELOG.txt .
As of September 2023, the EMBL-EBI’s FTP area hosts TAR
files for proteomes of 48 organisms, including model organ-
isms and WHO pathogens of interest (Supplementary Table
S1). The complete dataset is stored on Google Cloud Platform
(GCP) and is accessed through a file access API. The metadata
is indexed using Apache-Solr powering search API to facilitate
data accessibility and searchability.
The database provides access to over 214 million predicted
structures, although some sequences might be outdated com-
pared to UniProt due to less frequent data releases in the Al-
phaFold DB. Predictions of UniProt sequences are outputs
of a single model run. In contrast, Swiss-Prot / proteomes en-
tries represent the most confident prediction from runs of five
models trained with different random seeds. The following se-
quences are not covered in the database: (i) those that are less
than 16 amino acids, or (ii) > 2700 for SwissProt or proteome
sequences and 1280 for other UniProt sequences, or (iii) those
that contain non-standard amino acids, or (iv) are not in the
UniProt ‘one sequence per gene’ FASTA file, or (v) viral pro-
teins. These limitations are under discussion.
Data access
We facilitate access to the predicted structures and their as-
sociated confidence metrics from AlphaFold DB through four
different channels: FTP, Google Cloud Public Data, API, and
directly from the AlphaFold DB web page.
A subset of the AlphaFold DB can be accessed through
EMBL-EBI’s public FTP area at http:// ftp.ebi.ac.uk/ pub/
databases/ alphafold/ . The FTP area houses a comprehensive
README.txt file containing detailed information about all
available files. Predictions are archived according to versions,
with all versions, including the latest release, being accessi-
ble from folders within the FTP area. For example, the lat-
est archived files are available via http:// ftp.ebi.ac.uk/ pub/
databases/ alphafold/ latest/ . Additionally, supplementary files
such as sequences in FASTA format for all the predictions, a
CSV file listing UniProt accessions with predicted structures
and a CHANGELOG file highlighting version control are pro-
vided to help users. It is important to note that PAE data is
unavailable from the EMBL-EBI public FTP.
The full dataset, housing all predictions, is accessible from
Google Cloud Public Datasets under a CC-BY-4.0 license. This
dataset, approximately 23 TiB in size, is available at the fol-
lowing Google Cloud Storage Bucket: gs: // public-datasets-
deepmind-alphafold-v4. We suggest that most users down-
load only the subset of files relevant to their specific use case
to optimise resources. However, if a complete dataset is re-
quired for local processing, as might be the case in an academic
high-performance computing centre, it can be downloaded in
roughly 2.5 days using a 1 Gbps internet connection. Impor-
tantly, a Google account is necessary for the download.
The Alphafold DB API provides an efficient way for devel-
opers to programmatically access metadata associated with
all archived AlphaFold predictions. The API facilitates infor-
mation retrieval related to protein structures predicted by Al-
phaFold, such as URLs of model files (mmCIF, binaryCIF
and PDB), model quality metrics, and other valuable infor-
mation. All available API endpoints are keyed on UniProt ac-
cessions, and an interactive API documentation can be found
at https:// www.alphafold.ebi.ac.uk/ api-docs .
Finally, structure predictions can be directly accessed
through the AlphaFold DB website. The interface offers an
intuitive search functionality to quickly find and download
protein structure predictions and their corresponding confi-
dence metrics. Our commitment to enhancing user experience
and removing ambiguity has resulted in several improvements
to our user interface (UI) since its initial launch. One notable
advancement is the refined UI for search results, ensuring an
intuitive and user-friendly experience (Figure 2 ).
We have added filtering functionality to provide users with
a more tailored browsing experience. Users can now fil-
ter search results based on the status of the underlying se-
quence. This feature allows users to narrow their results to
only reviewed (Swiss-Prot) or unreviewed (TrEMBL) UniProt
Downloaded from https://academic.oup.com/nar/article/52/D1/D368/7337620 by guest on 29 April 2026
Nucleic Acids Research , 2024, Vol. 52, Database issue
D 371
Figure 2. Impro v e search results UI. T he impro v ed search UI includes more filtering options based on the underlying sequence of the protein str uct ures
and easier access to the most popular organisms.
accessions, enabling them to focus on predicted structures de-
rived from higher-quality, curated protein sequences.
An additional new option allows users to filter results based
on whether the protein is part of a UniProt reference proteome
dataset. This filter offers an additional indicator of confidence
in the quality of the sequence. Recognising the popularity of
specific organisms, we have also created a distinct list featur-
ing the most frequently searched-for species, allowing users
to swiftly find structures associated with popular species, like
Human, Mouse or Esc heric hia coli , enhancing our platform’s
overall user-friendliness and efficiency.
In response to frequent requests from the user community
since the AlphaFold DB launch, we have made significant
strides to implement a sequence-based similarity search fea-
ture. We have incorporated the Basic Local Alignment Search
Tool (BLAST) ( 21 ) into our Google Cloud Platform infrastruc-
ture to achieve this. This implementation allows our system
to swiftly compare user-provided protein sequences against
our database. We developed an API to send user protein se-
quences to the BLAST service and retrieve the responses. To
integrate these new search results into our internal search en-
gine, we built on the XJoin functionality in the BioSolr plu-
gin ( https:// github.com/ flaxsearch/ BioSolr ), further extending
it to suit our specific requirements. This functionality ensures
seamless integration of the BLAST search results and facili-
tates support for conventional filtering options.
We provide an intuitive user interface to present the se-
quence search results clearly and comprehensively (Figure 3 ).
The dedicated results page lists all similar proteins for which
we have predicted structures, and it offers capabilities for both
sorting and filtering, enhancing the ability to navigate the re-
sults efficiently.
With over 214 million predicted protein structures, the Al-
phaFold DB presents a colossal challenge in data analysis.
Tackling this challenge, Steinegger et al. recently unveiled the
Foldseek Cluster, a state-of-the-art structural-alignment-based
clustering algorithm designed for vast datasets ( 22 ). They ap-
plied this novel method to the AlphaFold DB archive, mak-
ing the derived structure similarity clusters available to the
research community. We integrated these structure similar-
ity clusters and continue to work on deploying a structure-
based similarity search into the AlphaFold DB. As part of
the initial rollout in our phased-release approach, tables have
been incorporated into the AFDB prediction pages, listing
AlphaFold predictions from the same similarity cluster as
the protein of interest (Figure 4 ). The clustering process was
two-pronged: the MMseqs2 tool first clustered 214 million
UniProtKB protein sequences from AlphaFold DB, trimming
it down to 52 million clusters based on defined sequence cri-
teria. A protein with the peak pLDDT score was elected as
each cluster’s representative. This set then underwent a sec-
ondary clustering via Foldseek, using specific structural de-
lineations, resulting in 18.8 million clusters. After dismiss-
ing sequences recognised as fragments, we finalised 2.30 mil-
lion robust clusters, each housing at least a pair of struc-
tures. We provide access to AlphaFold predictions from the
two main categories generated by the clustering process:
AFDB / Foldseek and AFDB50 / MMseqs. These results are also
accessible through the public API of AlphaFold DB.
The Predicted Aligned Error (PAE) plays a pivotal role as a
confidence metric to evaluate protein structures predicted by
AlphaFold. Since the initial launch of AlphaFold DB, we have
incorporated an interactive 2D heatmap visualisation on the
prediction pages. This visualisation tool allows users to focus
on specific regions and assess the confidence of AlphaFold’s
prediction regarding the regions’ relative positioning.
Understanding the importance of PAE in providing users
insights into the relative orientation of protein domains, we
have taken measures to enhance this feature further. We im-
proved the visualisation of non-consecutive regions and the
interactivity between the PAE viewer and the 3D molecular
graphics viewer, Mol* ( 23 ) (Figure 5 ). Now, when users select
Downloaded from https://academic.oup.com/nar/article/52/D1/D368/7337620 by guest on 29 April 2026
D 372
Nucleic Acids Research , 2024, Vol. 52, Database issue
Figure 3. Sequence similarity search results. We added support for performing sequence similarity searches in AlphaFold DB. The search results page
displa y s a list of predicted str uct ures with sequences similar to the user’s.
Figure 4. Str uct ure similarit y cluster members. Using dat a from AFDB Clusters, w e displa y lists of AlphaFold predictions str uct urally similar to a protein
of interest.
regions on the off-diagonal segment of the PAE heatmap, the
corresponding regions in the 3D view are highlighted. This
improvement not only bolsters the accessibility of the PAE
data but also empowers users to make informed interpreta-
tions concerning the overall conformation of the predicted
structure.
In addition to improving the interaction between Mol*
and the PAE viewer, we further expanded the capabilities of
our 3D molecular viewer to cater to more advanced analy-
sis in response to user feedback. Now, users can select in-
dividual atoms, residues, and complete chains, facilitating
a more comprehensive and focused exploration of molec-
ular structures. A practical application of this functional-
ity is measuring the distances between residue pairs, a fre-
quent action of in-depth investigation of protein structures
(Figure 6 ).
Downloaded from https://academic.oup.com/nar/article/52/D1/D368/7337620 by guest on 29 April 2026
Nucleic Acids Research , 2024, Vol. 52, Database issue
D 373
Figure 5. Impro v ed support f or highlighting non-consecutiv e regions. T he ne w v ersion of the interactiv e PAE vie w er mak es it easier to distinguish
between highlighted non-consecutive regions when assessing their relative positions’ confidence, as shown for AlphaFold DB accession
https:// alphafold.ebi.ac.uk/ entry/ Q7RTU9 .
Figure 6. Impro v ed customisation in Mol*. Enhanced customisation options in Mol* allo w users to perf orm popular actions such as measuring
distances or changing the rendering style.
Downloaded from https://academic.oup.com/nar/article/52/D1/D368/7337620 by guest on 29 April 2026
D 374
Nucleic Acids Research , 2024, Vol. 52, Database issue
Conclusions and outlook
Accessing hundreds of millions of predicted protein structures
housed within data resources such as the AlphaFold Protein
Structure Database marks a significant leap for structural bi-
ology and has impacted diverse fields within the life sciences.
The impact of the predicted structures is further enhanced by
their seamless integration into several core data resources, in-
cluding PDBe-KB ( 17 ), UniProt ( 6 ), Ensembl ( 16 ), InterPro
( 15 ), Genecards ( 24 ) and MobiDB ( 25 ), among others. How-
ever, while the field has made significant strides, there are dis-
cernible gaps in both data representation and functionality
that we are committed to addressing.
There are intriguing frontiers for data enhancement, in-
cluding adding structures for isoforms and targeted datasets
for multimeric predicted protein structures. In parallel with
these anticipated data updates, we are also working towards
enriching the predicted structures with domain annotations
( 26 ), integrating small molecules through AlphaFill ( 27 ), ref-
erencing cross-linking data, and devising dedicated pages tai-
lored for fragments. These planned improvements are based
on our ongoing interactions with the scientific community,
whose feedback and insights highlight areas where the most
impactful changes could be made. We invite all users to share
their suggestions through the AlphaFold DB helpdesk (afdb-
help@ebi.ac.uk).
While the availability of millions of predicted protein struc-
tures promises valuable insights into molecular biology, they
might be behind a barrier for many researchers who may lack
familiarity with handling macromolecular structure data and
may not sufficiently understand the strengths and limitations
inherent in predicted structures. To address these challenges
and make protein structure data more accessible, we now fo-
cus on providing a comprehensive training platform, enabling
the broader scientific community to use structural data more
efficiently.
Arguably, we have entered a new era in structural biology,
when the abundance of available predicted protein structure
data enables researchers to probe an unprecedented range of
biological questions. As stewards of AlphaFold DB, we are
committed to bolstering its accessibility, hoping to amplify its
transformative impact on science and society.
Data availability
V ersioned T AR files for 48 model organisms and pathogens
and metadata files are available from the public FTP area of
EMBL-EBI at http:// ftp.ebi.ac.uk/ pub/ databases/ alphafold/ .
The documentation of the public API endpoints of AlphaFold
DB is available at https:// www.alphafold.ebi.ac.uk/ api-docs .
A guide on accessing all the Google Cloud Public Dataset data
is available from our download page https://www.alphafold.
ebi.ac.uk/download .
Supplementary data
Supplementary Data are available at NAR Online.
A c kno wledg ements
We acknowledge the researchers who consistently contribute
to depositing structures, enriching the foundation of structural
biology and the primary databases, notably UniProt, MGnify
and PDB, for providing open data that is indispensable in
training structure prediction tools. Lastly, we thank the PDBe
and Google DeepMind teams for their rigorous evaluation of
new features; their collective contributions have been invalu-
able to our work.
Funding
Google DeepMind funds the AlphaFold Protein Struc-
ture Database. Funding for open access charge: Google
DeepMind funds. M.S. acknowledges support from
the National Research Foundation of Korea (grants
2019R1A6A1A10073437,
2020M3A9G7103933,
2021R1C1C102065, and 2021M3A9I4021220), the Sam-
sung DS Research Fund, and the Creative-Pioneering Re-
searchers Program through Seoul National University. M.M.
acknowledges support by the National Research Foundation
of Korea (grant RS-2023-00250470).
Conflict of interest statement
None declared.
References
1. Jumper, J. , Evans, R. , Pritzel, A. , Green, T. , Figurnov, M. ,
Ronneberger, O. , Tunyasuvunakool, K. , Bates, R. , Žídek, A. ,
Potapenko, A. , et al. (2021) Highly accurate protein structure
prediction with AlphaFold. Nature , 596 , 583–589.
2. Baek, M. , DiMaio, F. , Anishchenko, I. , Dauparas, J. , Ovchinnikov, S. ,
Lee, G.R. , Wang, J. , Cong, Q. , Kinch, L.N. , Schaeffer, R.D. , et al.
(2021) Accurate prediction of protein structures and interactions
using a three-track neural network. Science , 373 , 871–876.
3. Ahdritz, G. , Bouatta, N. , Floristean, C. , Kadyan, S. , Xia, Q. ,
Gerecke, W. , O’Donnell, T.J. , Berenberg, D. , Fisk, I. , Zanichelli, N. ,
et al. (2022) OpenFold: retraining AlphaFold2 yields new insights
into its learning mechanisms and capacity for generalization
Bioinformatics. bioRxiv doi:
https:// doi.org/ 10.1101/ 2022.11.20.517210 , 22 November 2022,
preprint: not peer reviewed.
4. Kryshtafovych, A. , Schwede, T. , Topf, M. , Fidelis, K. and Moult, J.
(2021) Critical assessment of methods of protein structure
prediction (CASP)-Round XIV. Proteins , 89 , 1607–1617.
5. Velankar, S. , Burley, S.K. , Kurisu, G. , Hoch, J.C. and Markley, J.L.
(2021) The Protein Data Bank Archive. Methods Mol. Biol.
Clifton NJ , 2305 , 3–21.
6. U.P. Consortium. (2023) UniProt: the Universal Protein
Knowledgebase in 2023. Nucleic Acids Res. , 51 , D523–D531.
7. Varadi, M. , Bordin, N. , Orengo, C. and Velankar, S. (2023) The
opportunities and challenges posed by the new generation of deep
learning-based protein structure predictors. Curr. Opin. Struct.
Biol., 79 , 102543.
8. Bordin, N. , Dallago, C. , Heinzinger, M. , Kim, S. , Littmann, M. ,
Rauer, C. , Steinegger, M. , Rost, B. and Orengo, C. (2023) Novel
machine learning approaches revolutionize protein knowledge.
Trends Biochem. Sci., 48 , 345–359.
9. Mosalaganti, S. , Obarska-Kosinska, A. , Siggel, M. , Taniguchi, R. ,
Turo ˇnová, B. , Zimmerli, C.E. , Buczak, K. , Schmidt, F.H. ,
Margiotta, E. , Mackmull, M.-T. , et al. (2022) AI-based structure
prediction empowers integrative structural analysis of human
nuclear pores. Science , 376 , eabm9506.
10. Goverde, C.A. , Wolf, B. , Khakzad, H. , Rosset, S. and Correia, B.E.
(2023) De novo protein design by inversion of the AlphaFold
structure prediction network. Protein Sci. Publ. Protein Soc., 32 ,
e4653.
11. Bordin, N. , Sillitoe, I. , Nallapareddy, V. , Rauer, C. , Lam, S.D. ,
Waman,V .P ., Sen,N., Heinzinger,M., Littmann,M., Kim,S., et al.
Downloaded from https://academic.oup.com/nar/article/52/D1/D368/7337620 by guest on 29 April 2026
Nucleic Acids Research , 2024, Vol. 52, Database issue
D 375
(2023) AlphaFold2 reveals commonalities and novelties in protein
structure space for 21 model organisms. Commun. Biol., 6 , 160.
12. Fontana, P. , Dong, Y. , Pi, X. , Tong, A.B. , Hecksel, C.W. , Wang, L. ,
Fu, T.-M. , Bustamante, C. and Wu, H. (2022) Structure of
cytoplasmic ring of nuclear pore complex by integrative cryo-EM
and AlphaFold. Science , 376 , eabm9326.
13. Nussinov, R. , Zhang, M. , Liu, Y. and Jang, H. (2023) AlphaFold,
allosteric, and orthosteric drug discovery: ways forward. Drug
Discov. Today , 28 , 103551.
14. Lin, Z. , Akin, H. , Rao, R. , Hie, B. , Zhu, Z. , Lu, W. , Smetanin, N. ,
Verkuil, R. , Kabeli, O. , Shmueli, Y. , et al. (2022) Evolutionary-scale
prediction of atomic level protein structure with a language model.
bioRxiv doi: https:// doi.org/ 10.1101/ 2022.07.20.500902 , 31
October 2022, preprint: not peer reviewed.
15. Paysan-Lafosse, T. , Blum, M. , Chuguransky, S. , Grego, T. , Pinto, B.L. ,
Salazar, G.A. , Bileschi, M.L. , Bork, P. , Bridge, A. , Colwell, L. , et al.
(2023) InterPro in 2022. Nucleic Acids Res. , 51 , D418–D427.
16. Cunningham, F. , Allen, J.E. , Allen, J. , Alvarez-Jarreta, J. ,
Amode, M.R. , Armean, I.M. , Austine-Orimoloye, O. , Azov, A.G. ,
Barnes, I. , Bennett, R. , et al. (2022) Ensembl 2022. Nucleic Acids
Res., 50 , D988–D995.
17. consortium,P.D.B.-K.B. (2022) PDBe-KB: collaboratively defining
the biological context of structural data. Nucleic Acids Res., 50 ,
D534–D542.
18. Vallat, B. , Tauriello, G. , Bienert, S. , Haas, J. , Webb, B.M. , Žídek, A. ,
Zheng, W. , Peisach, E. , Piehl, D.W. , Anischanka, I. , et al. (2023)
ModelCIF: an Extension of PDBx / mmCIF Data Representation
for Computed Structure Models. J. Mol. Biol., 435 , 168021.
19. Mariani, V. , Biasini, M. , Barbato, A. and Schwede, T. (2013) lDDT: a
local superposition-free score for comparing protein structures and
models using distance difference tests. Bioinformatics , 29 ,
2722–2728.
20. Morales, J. , Pujar, S. , Loveland, J.E. , Astashyn, A. , Bennett, R. ,
Berry, A. , Cox, E. , Davidson, C. , Ermolaeva, O. , Farrell, C.M. , et al.
(2022) A joint NCBI and EMBL-EBI transcript set for clinical
genomics and research. Nature , 604 , 310–315.
21. Camacho, C. , Coulouris, G. , Avagyan, V. , Ma, N. , Papadopoulos, J. ,
Bealer, K. and Madden, T.L. (2009) BLAST+: architecture and
applications. BMC Bioinf., 10 , 421.
22. Barrio-Hernandez, I. , Yeo, J. , Jänes, J. , Mirdita, M. , Gilchrist, C.L.M. ,
Wein, T. , Varadi, M. , Velankar, S. , Beltrao, P. and Steinegger, M.
(2023) Clustering-predicted structures at the scale of the known
protein universe. Nature , 622 , 637–645.
23. Sehnal, D. , Bittrich, S. , Deshpande, M. , Svobodová, R. , Berka, K. ,
Bazgier, V. , Velankar, S. , Burley, S.K. , Ko ˇca, J. and Rose, A.S. (2021)
Mol* Viewer: modern web app for 3D visualization and analysis
of large biomolecular structures. Nucleic Acids Res., 49 ,
W431–W437.
24. Stelzer, G. , Rosen, N. , Plaschkes, I. , Zimmerman, S. , Twik, M. ,
Fishilevich, S. , Stein, T.I. , Nudel, R. , Lieder, I. , Mazor, Y. , et al. (2016)
The GeneCards Suite: from gene data mining to disease genome
sequence analyses. Curr. Protoc. Bioinforma., 54 , 1.30.1–1.30.33.
25. Piovesan, D. , Necci, M. , Escobedo, N. , Monzon, A.M. , Hatos, A. ,
Mi ˇceti ´c, I. , Quaglia, F. , Paladin, L. , Ramasamy, P. , Dosztányi, Z. , et al.
(2021) MobiDB: intrinsically disordered proteins in 2021. Nucleic
Acids Res., 49 , D361–D367.
26. Wells, J. , Hawkins-Hooker, A. , Bordin, N. , Paige, B. and Orengo, C.
(2023) Chainsaw: protein domain segmentation with fully
convolutional neural networks Molecular Biology. bioRxiv doi:
https:// doi.org/ 10.1101/ 2023.07.19.549732 , 19 July 2023,
preprint: not peer reviewed.
27. Hekkelman, M.L. , de Vries, I. , Joosten, R.P. and Perrakis, A. (2023)
AlphaFill: enriching AlphaFold models with ligands and cofactors.
Nat. Methods , 20 , 205–213.
Received: September 29, 2023. Revised: October 13, 2023. Editorial Decision: October 16, 2023. Accepted: October 18, 2023
© The Author(s) 2023. Published by Oxford University Press on behalf of Nucleic Acids Research.
This is an Open Access article distributed under the terms of the Creative Commons Attribution License (http: // creativecommons.org / licenses / by / 4.0 / ), which permits unrestricted reuse,
distribution, and reproduction in any medium, provided the original work is properly cited.
Downloaded from https://academic.oup.com/nar/article/52/D1/D368/7337620 by guest on 29 April 2026
\end{document}
